\documentclass[12pt,a4paper]{article}
\usepackage[utf8]{inputenc}
\usepackage[T1]{fontenc}
\usepackage{geometry}
\geometry{margin=1in}
\usepackage{enumitem}
\usepackage{listings}
\usepackage{xcolor}
\usepackage{hyperref}
\usepackage{titlesec}
\usepackage{fancyhdr}

% Code listing style
\lstset{
    basicstyle=\ttfamily\small,
    breaklines=true,
    frame=single,
    backgroundcolor=\color{gray!10},
    keywordstyle=\color{blue},
    commentstyle=\color{green!50!black},
    stringstyle=\color{red},
    showstringspaces=false,
    tabsize=4
}

\pagestyle{fancy}
\fancyhf{}
\fancyhead[L]{100 Practice Questions of String in Python}
\fancyhead[R]{CS-IP Learning Hub}
\fancyfoot[C]{\thepage}
\renewcommand{\headrulewidth}{0.4pt}

\title{\textbf{100 Important Practice Questions of String in Python}}
\author{Compiled from CS-IP-Learning-Hub\\
\url{https://csiplearninghub.com/practice-questions-of-string-in-python/}}
\date{}

\begin{document}

\maketitle
\tableofcontents
\newpage

\section{Practice Questions of String in Python -- Test 1}

\begin{enumerate}[label=\textbf{Q\arabic*.}]

\item What is String in Python?

\item Is there any difference in `a' or ``a'' in python?

\item Is there any difference between 1 or `1' in python?

\item Python treats single quotes same as double quotes. (T/F)

\item A string with zero character is called \underline{\hspace{3cm}} string.

\item Python does not support a character type. (T/F)

\item Write a code to store following String in a variable named `str'. \\
\textit{This is Amit's Blog}

\item Write the output of the following code:
\begin{lstlisting}[language=Python]
str1 = "Welcome to my blog"
str2 = "Welcome to my \n Blog"
print(str1)
print(str2)
\end{lstlisting}

\item Write the output of the following.
\begin{lstlisting}[language=Python]
str1 = "Welcome \tto my Blog"
str2 = "Welcome to\n my \tBlog"
print(str1)
print(str2)
\end{lstlisting}

\item Write the output of the following.
\begin{lstlisting}[language=Python]
str1 = """ Welcome to my
blog.
This is for
Class X"""
print(str1)
\end{lstlisting}

\end{enumerate}

\section{Practice Questions of String in Python -- Test 2}

\begin{enumerate}[label=\textbf{Q\arabic*.}]

\item Write the output of following code- 

\textbf{a)}
\begin{lstlisting}[language=Python]
str="hello"
print(str[:3])
\end{lstlisting}

\textbf{b)}
\begin{lstlisting}[language=Python]
str='My Blog'
a=' '
for i in range(len(str)):
    a+=str[i]
print(a)
\end{lstlisting}

\textbf{c)}
\begin{lstlisting}[language=Python]
str='MyBLog'
a=' '
for i in range(len(str)):
    print(i * str[i])
\end{lstlisting}

\textbf{d)}
\begin{lstlisting}[language=Python]
s='My'
s1='Blog'
s2=s[:1]+s1[len(s1)-1:]
print(s2)
\end{lstlisting}

\textbf{e)} \texttt{print("My"+'Blog' * 2)}

\textbf{f)} \texttt{print("My" *3 + "Blog" +'7')}

\textbf{g)}
\begin{lstlisting}[language=Python]
for i in range(2,7,2):
    print(i * '2')
\end{lstlisting}

\textbf{h)}
\begin{lstlisting}[language=Python]
for i in range(3,12, 2):
    print("a".upper())
\end{lstlisting}

\textbf{i)}
\begin{lstlisting}[language=Python]
for i in range(3,12,13):
    print("a".upper)
\end{lstlisting}

\textbf{j)}
\begin{lstlisting}[language=Python]
s='Welcome to My Site https://csiplearninghub.com '
print(s.find('come'))
print(s.find('o'))
print(s.find('o', 10, 20))
print(s.find('o', 5, 10))
\end{lstlisting}

\end{enumerate}

\section{Practice Questions of String in Python -- Test 3}

\begin{enumerate}[label=\textbf{Q\arabic*.}]

\item Write a code to create empty string `str1'

\item What do you mean by traversing a string?

\item What is the index value of first element of a string?

\item What is the index value of last element of a string?

\item If the length of the string is 10 then what would be the positive index value of last element?

\item If the length of string is 9, what would be the index value of middle element?

\item Index value of a string can be in float. (T/F)

\item What type of error is returned by following statement, if the length of string `str1' is 10.
\begin{lstlisting}[language=Python]
print(str1[13])
\end{lstlisting}

\item Write the output of the following:
\begin{lstlisting}[language=Python]
str1 = "Welcome to my Blog"
a. print(str1[-1])
b. print(str1[9])
\end{lstlisting}

\item Write a code to assign a string ``Hello World'' to a string variable named ``str1''.

\end{enumerate}

\section{Practice Questions of String in Python -- Test 4}

\begin{enumerate}[label=\textbf{Q\arabic*.}]

\item Write a program to display each character of the following string in separate line using `for' loop. \\
\texttt{str1 = Welcome to My Blog}

\item Write a program to display each character of the following string in separate line using `while' loop. \\
\texttt{str1 = Welcome to My Blog}

\item Index value of string ``str1'' varies from 0 to len(str1)-1. (T/F)

\item Write the positive and negative index value of `B' in the following string. \\
\texttt{"Welcome to my Blog"}

\item What do you mean by concatenation of string?

\item Which of the following is an example of concatenation?       
\begin{enumerate}[label=\alph*.]
\item 6 + 3
\item `6' + `3'
\item `a' + `b' + `c'
\end{enumerate}

\item Write a program to count the length of string without using inbuilt function.

\item Write the output of the following statement.
\begin{lstlisting}[language=Python]
str1 = "Welcome to my Blog"
for i in str1:
    print(i)
    print(i, end =' ')
    print(i, end = '#')
\end{lstlisting}

\item Write the output of the following statement.
\begin{lstlisting}[language=Python]
str1 = "Amit"
for i in str1:
    print(i)
    print(i, end =' ')
\end{lstlisting}

\item Write the output of the following statement.
\begin{lstlisting}[language=Python]
str1 = "Welcome to my Blog"
for i in str1:
    print(i)
    print(i, end =' ')
print(i, end = '#')
\end{lstlisting}

\end{enumerate}

\section{Practice Questions of String in Python -- Test 5}

\begin{enumerate}[label=\textbf{Q\arabic*.}]

\item Match the Following

\begin{center}
\begin{tabular}{|l|l|}
\hline
\textbf{Column-1} & \textbf{Column-2} \\
\hline
Slicing & string[range] + `x' \\
Concatenation & string[: -- 1] \\
Repetition & string[range] \\
Membership & string1 + string2 \\
Reverse & in, not in \\
\hline
\end{tabular}
\end{center}

\item Write the code to join string `str1' and `str2' and store the result in variable `str3'

\item The \underline{\hspace{3cm}} operator join the string.

\item What type of error is shown by following code? \\
\texttt{>>> 2 + '3'}

\item Which operator is used to replicate string?

\item Write the output of the following. \\
\texttt{>>> 'A' * 3}

\item Write the output of the following:
\begin{lstlisting}[language=Python]
for i in range(2,5,2):
    print(i * '@')
\end{lstlisting}

\item `in' or `not in' are the \underline{\hspace{3cm}} operator

\item Write the output of the following
\begin{lstlisting}[language=Python]
>>> 'h' in 'Hindi'
>>>'i' in 'Hindi'
\end{lstlisting}

\item Write the output of the following
\begin{lstlisting}[language=Python]
for i in "Amit":
    print(i in "Amit")
\end{lstlisting}

\end{enumerate}

\section{Practice Questions of String in Python -- Test 6}

\begin{enumerate}[label=\textbf{Q\arabic*.}]

\item Write the output of the following
\begin{lstlisting}[language=Python]
>>>"string" in "substring"
>>>"string" not in "substring"
\end{lstlisting}

\item Write the output of the following:
\begin{lstlisting}[language=Python]
>>>'tie' == 'ties'
>>>"Amit" != "Amitabh"
>>>"amit" > "Amitabh"
\end{lstlisting}

\item Write the full form of ASCII

\item Write the ASCII value of `A' and `a'.

\item String Slicing is used to fetch substring from a string. (T/F)

\item Write the output of the following.
\begin{lstlisting}[language=Python]
S = "Welcome to my Blog"
print(S[2 : 3])
print(S[2 : 10])
print(S[-2 : ])
print(S[-10 : -2 : 2])
\end{lstlisting}

\item Strings are mutable. (T/F)

\item Write the output of the following
\begin{lstlisting}[language=Python]
>>>str1 = "Anita"
>>>str1[1]  = 'm'
>>>print(str1)
\end{lstlisting}

\item Name any two built in function associated with string.

\item Which function of string return the numerical value?

\end{enumerate}

\section{Practice Questions of String in Python -- Test 7}

\begin{enumerate}[label=\textbf{Q\arabic*.}]

\item What is the value returned by find() function for an unsuccessful search?

\item isalpha() returns True if the string contains only alphabet. (T/F)

\item Write the output of the following
\begin{lstlisting}[language=Python]
str1 = "Amit"
str2 = "My Blog"
str3 = "#blog"
str4 = "My 1st Blog"
print(str1.isalpha())
print(str2.isalpha())
print(str3.isalpha())
print(str4.isalpha())
\end{lstlisting}
Specify the reason if any of the above print statement return False.

\item Write a program to accept string and display total number of alphabets.

\item Write a program to accept a string and display the sum of the digits, if any present in string. \\
for example: \\
input string : My position is 1st and my friend come on 4th \\
output : 5

\item Write the output of the following:
\begin{lstlisting}[language=Python]
>>>a = 123
>>>b = "123"
>>>b.isdigit()
>>>a.isdigit()
\end{lstlisting}

\item What is the difference between lower() and islower()?

\item Write a program to accept a string and convert it into lowercase.

\item Write a program to count the number of lowercase and uppercase character in a string accepted from the user.

\item Write the output of the following:
\begin{lstlisting}[language=Python]
>>>s = "My blog"
>>>s.upper()
>>>s.lower()
>>>s.islower()
>>>s.isupper()
>>>s.isalpha()
>>>s.isdigit()
\end{lstlisting}

\end{enumerate}

\section{Practice Questions of String in Python -- Test 8}

\begin{enumerate}[label=\textbf{Q\arabic*.}]

\item What is lstrip() function in String?

\item Write the output of the following. (\# represents the spaces)
\begin{lstlisting}[language=Python]
>>>str1 = "Welcome to my Blog"
>>>print(len(str1))
>>>str2 = "###Welcome to my Blog"
>>>print(len(str2))
>>>print(len(str2.lstrip()))
\end{lstlisting}

\item Write the output of the following. (\# represents the spaces)
\begin{lstlisting}[language=Python]
>>>str2 = "###Welcome to my Blog####"
>>>print(len(str2))
>>>print(len(str2.strip()))
>>>print(len(str2.rstrip()))
\end{lstlisting}

\item Write the output of the following.
\begin{lstlisting}[language=Python]
>>>str1 = "Welcome to my Blog"
>>>print(str1.rstrip('og'))
>>>print(str1.lstrip('We'))
>>>print(str1.strip('Welog'))
\end{lstlisting}

\item Write the output of the following.
\begin{lstlisting}[language=Python]
>>>print(len(str1.rstrip('og')))
>>>print(len(str1.lstrip('We')))
>>>print(len(str1.strip('Welog')))
\end{lstlisting}

\item isspace() returns True if the string contain only spaces (T/F)

\item Write the output of the following
\begin{lstlisting}[language=Python]
>>>s1 = "   "
>>>s2 = "   Amit"
>>>print(s1.isspace())
>>>print(s2.isspace())
\end{lstlisting}

\item Define the following function with example
\begin{enumerate}[label=\alph*.]
\item istitle()
\item swapcase()
\end{enumerate}

\end{enumerate}

\section{Practice Questions of String in Python -- Test 9}

\begin{enumerate}[label=\textbf{Q\arabic*.}]

\item Write the output of the following:
\begin{lstlisting}[language=Python]
str1 = "Welcome to my Blog"
a. print(len(str1))
b. print(capitalize(str1))
\end{lstlisting}

\item Write the output of the following.
\begin{lstlisting}[language=Python]
>>>str1 = "Welcome to my Blog"
>>>x = str1.split()
>>>print(x)
\end{lstlisting}

\item Write a program to accept a string and display each word and it's length.

\item Write a program to accept a string and display string with capital letter of each word. \\
for example, if input string is : welcome to my blog \\
output string : Welcome To My Blog

\item What is split() function in String?

\item Write the output of the following:
\begin{lstlisting}[language=Python]
a = "Mummy?Papa?Brother?Sister?Uncle"
print(a.split())
print(a.split('?'))
print(a.split('?',1))
print(a.split('?',3))
print(a.split('?',10))
print(a.split('?',-1))
\end{lstlisting}

\item Write a program to replace all the word `do' with `done' in the following string. \\
\texttt{str1 = "I do you do and we all will do"}

\item Write the output of the following.
\begin{lstlisting}[language=Python]
str1 = "I went to Auli"
print(str1.replace("Auli", "Leh"))
print(str1)
\end{lstlisting}

\item What is the purpose of find() function in string?

\item Write the output of the following:
\begin{lstlisting}[language=Python]
str1 = "Welcome to my Blog"
print(str1.find('o'))
print(str1.find('o',3))
print(str1.find('o',7))
print(str1.find('o',7,10))
\end{lstlisting}

\end{enumerate}

\section{Practice Questions of String in Python -- Test 10}

\begin{enumerate}[label=\textbf{Q\arabic*.}]

\item Accept a string and display in reverse order.

\item Write a program to accept a string and display 20 times.

\item Write a program to accept a string in python and display the entire string in upper case.

\item Write a program to accept a string and display last three characters of the string.

\item Write a program to accept a string in python and display the entire string in lower case.

\item Accept a string and display the entire string with first and last character in upper case.

\item Write a program to accept a string and display first three characters of the string.

\item Write a function in python which accept a string as argument and display total number of vowels.

\item Write a function in python which accept a string as argument and display total number of lower case characters.

\item Write a function in python which accept a string as argument and display total number of digits.

\end{enumerate}

\vspace{1cm}
\begin{center}
\textit{Source: CS-IP-Learning-Hub}\\
\url{https://csiplearninghub.com/practice-questions-of-string-in-python/}
\end{center}

\end{document}
